\section{Additional Data for some Specific Zones} \label{sec2.8}

\medskip

\textbf{General remarks}

\medskip

Input data in this block permit explicit specification of plasma
parameters $T_e$ , $T_i$ , $D_i$, $V_x$, $V_y$, $V_z$ and of zone
volumes VOL in selected cells ICELL. Furthermore, information may
be given to the geometrical block of EIRENE that some ``additional
surfaces" are ``invisible" for a particle located in cell ICELL
and, therefore, possible crossings need not be checked for
advancing this particle at the next step. (Intelligent use of this
option can lead to a considerable speeding up of the code, less
intelligent use will lead to dramatic errors.) All this
information will be included in the EIRENE arrays after the
initialization phase (Subroutines INPUT, PLASMA, GRID, VOLUME) and
before setting the derived profiles ($D_e$, flux-surface labelling
grids, \dots)

\medskip

\textbf{The Input Block}

\medskip

\textsl{*** 8. Data for some specific zones}
\begin{lstlisting}
NZADD
DO 81 IZADD=1,NZADD
*   comment (optional card)
  INI INE
\end{lstlisting}
then read an arbitrary number of cards
each starting with either T, D, V, M, VL or CH3 in
any order

\begin{lstlisting}
81  CONTINUE
\end{lstlisting}

The following specifications are applied to each cell with
cell number ICELL in the range  $INI \leq ICELL \leq INE$.
\begin{list}{}{}
\item[\textbf{``CH3-cards"}] (format: 3A,69A)

arbitrary number of strings $\pm$ n/m, separated by blanks. n and
m must be integer variables 1 $\le$ n,m $\le$ NLIMI. By default
EIRENE assumes that any additional surface is ``visible" from any
cell of the computational box. A string $n/m$ (or $+n/m$) has the
effect that during particle tracing possible crossings of
additional surfaces number $n$ to number $m$ are not checked whenever
a track starts in cell ICELL. By $-n/m$ these surfaces are activated
again (in principle this option is not needed due to default
setting). 

A particle history can be forced to stop and then to
restart in a cell e.g.\ by making appropriate use of the ILIIN = -1
option for additional- or non default standard surfaces.
\item[\textbf{``T-cards"}] T, IDION, TEADD, TIADD(IDION)

format: 1A,1I6,6X,2E12.4

plasma temperatures $T_e $ and $T_i $ in cell
ICELL are reset to:

TEIN(ICELL)=TEADD

TIIN(IDION,ICELL)=TIADD(IDION)
\item[\textbf{``D-cards"}] D, IDPLS,
DIADD(IDPLS)

format: 1A,1I6,6X,1E12.4

plasma ion density $D_i $ for species IDPLS in cell
ICELL is reset to:

DIIN(IDPLS,ICELL)=DIADD(IDPLS) \item[\textbf{``V-cards"}] V, IDPLS,
VXADD(IDPLS), VYADD(IDPLS), VZADD(IDPLS)

format: 1A,1I6,6X,3E12.4

plasma drift velocity in x direction
$V_x $ for species IDPLS in cell
ICELL is reset to:

VXIN(IDPLS,ICELL)=VXADD(IDPLS)

VYIN,VZIN: likewise, for the drift velocities in y and z direction
respectively.
\item[\textbf{``M-cards"}] M, IDPLS, MXADD(IDPLS),
MYADD(IDPLS), MZADD(IDPLS)

same as ``V-cards", but drift velocity in cell ICELL is given in
Mach number units:

$VXIN = MXADD  \cdot
\frac{\sqrt{T_e + T_i}}   {m_i} $ ,

$m_i $ being the mass of bulk ions of species IDPLS.

VYIN,VZIN are computed in the same way from MYADD and MZADD
respectively.
\item[\textbf{``VL-cards"}] VL, VLADD

format: 2A,1E12.4

The zone volume in cell ICELL is explicitly set to:

VOL(ICELL)=VLADD
\end{list}
